\section{Independent replication on NIH radiographs}
\label{app:seed-results}

\begin{table*}[pos=tp]
\centering
\setlength{\tabcolsep}{3pt}
% terracotta = positive / high, blue = negative / low, sage = magnitude of Read and Ans; darker = larger magnitude
\definecolor{cfSage1}{HTML}{ECF3EF}
\definecolor{cfSage2}{HTML}{D5E5DB}
\definecolor{cfSage3}{HTML}{B9D4C3}
\definecolor{cfSage4}{HTML}{9EC3AC}
\definecolor{cfSlate1}{HTML}{EAF0F6}
\definecolor{cfSlate2}{HTML}{D0DDEA}
\definecolor{cfSlate3}{HTML}{B1C7DD}
\definecolor{cfSlate4}{HTML}{94B2D0}
\definecolor{cfTerra1}{HTML}{F7ECEA}
\definecolor{cfTerra2}{HTML}{EED5D1}
\definecolor{cfTerra3}{HTML}{E2B8B2}
\definecolor{cfTerra4}{HTML}{D79E95}
\definecolor{cfOat1}{HTML}{F4EFE6}
\definecolor{cfOat2}{HTML}{E8DDC8}
\definecolor{cfGrey}{HTML}{E4E1DC}

\caption{\textbf{Independent replication on NIH radiographs.} Grades and steering effects for Qwen2.5-VL-7B. The random 95th percentile (p95) and absolute sham effect are the steering references. Blue marks a negative row margin.}
\label{tab:cf-seed}
\fitwidth{%
\cftabsize
\begin{tabular}{lrrrrrrrll}
\toprule
Concept & \multicolumn{1}{c}{$S$} & \multicolumn{1}{c}{AUROC} & \multicolumn{1}{c}{$W_{q,q}$} & \multicolumn{1}{c}{p95} & \multicolumn{1}{c}{$|$sham$|$} & \multicolumn{1}{c}{rank} & \multicolumn{1}{c}{$O_q$} & Competitor & Verdict\\
\midrule
Effusion & 0.102 & 0.593 & +0.190 & 0.137 & 0.017 & 2/120 & \cellcolor{cfSlate2}-0.071 & Nodule & competitor \\
Atelectasis & 0.062 & 0.566 & +0.101 & 0.297 & 0.096 & 55/120 & \cellcolor{cfSlate3}-0.226 & Nodule & competitor \\
Pneumothorax & 0.152 & 0.508 & +0.089 & 0.320 & 0.113 & 52/120 & \cellcolor{cfSlate3}-0.102 & Effusion & competitor \\
Cardiomegaly & 0.205 & 0.742 & +0.085 & 0.311 & 0.158 & 33/120 & \cellcolor{cfSlate3}-0.219 & Atelectasis & competitor \\
Mass & 0.021 & 0.678 & +0.359 & 0.372 & 0.092 & 10/120 & \cellcolor{cfSlate3}-0.169 & Effusion & competitor \\
Nodule & -0.041 & 0.642 & +0.118 & 0.158 & 0.055 & 15/120 & \cellcolor{cfSlate2}-0.072 & Effusion & competitor \\
\bottomrule
\end{tabular}}
\end{table*}


\subsection{A non-generic write effect reproduces}
\label{sec:qwen-replication}

The first finding is a successful intervention. Qwen Effusion reproduces beyond random and sham controls, and only the clinical alternatives tested afterwards change its attribution.

Controlled decodability is evaluated at the final visual block each connector consumes: the penultimate CLIP encoder layer for LLaVA and the last of Qwen's 32 visual blocks. The control AUROC is the mean over the 20 recurring-type random-label maps.

Qwen Effusion first satisfies the reader test. Its probe reaches AUROC \cfSeedDecQwenEffAuroc{} against a control AUROC of \cfSeedDecQwenEffCtrl{}, giving a controlled selectivity of \cfSeedDecQwenEffSel{}. LLaVA-1.5-7B reads both of its concepts at its selected vision layer, against a control AUROC of \cfSeedDecLlavaEffCtrl{}. Effusion reaches AUROC \cfSeedDecLlavaEffAuroc{} with selectivity \cfSeedDecLlavaEffSel{}, and Edema reaches AUROC \cfSeedDecLlavaEdemaAuroc{} with selectivity \cfSeedDecLlavaEdemaSel{}. For the original-cohort interventions, we take the maximum concept-consistent effect over six fixed nonzero strengths. In the original 200-image cohort, Qwen's maximum concept-consistent change is 0.2343 at $\alpha=+0.25$. This exceeds the random-direction 95th percentile of 0.0963 and the maximum absolute sham effect of 0.0852, a strong case for the write direction.

The response at the locked steering strength $\alpha=+0.25$ reproduces on the patient-disjoint 400-patient cohort. Effusion changes mean $P(\mathrm{yes})$ by 0.1945, again above the random-direction 95th percentile of 0.0878 and absolute sham effect of 0.0157. None of these patients appears in the original intervention cohort, and the steering strength was not reselected. The result supports a large, reproducible, non-generic effect within ChestX-ray14.


\subsection{Competing directions on the effusion question}
\label{sec:specificity-results}

All five alternative directions increased the effusion answer. The nodule direction produced the largest increase, giving the effusion answer an ownership margin of $O_{\mathrm{Effusion}}=\cfSeedPriorOFull$. The effusion direction met the random and sham references. The same direction was outperformed by the nodule direction on a separate test cohort (Table~\ref{tab:cf-seed}).
